% ============================================================
%  TUTORIAL COMPLETO: API REST con Node.js
%  Gestión de Cultivos -- AgroAPI Ecuador
%  Compilar con: pdflatex tutorial_api.tex
% ============================================================
\documentclass[12pt, a4paper]{article}

% -- Codificación y fuentes -----------------------------------
\usepackage[utf8]{inputenc}
\usepackage[T1]{fontenc}
\usepackage[spanish]{babel}
\usepackage{lmodern}

% -- Geometría de página --------------------------------------
\usepackage[
  top=2.5cm, bottom=2.5cm,
  left=2.8cm, right=2.8cm
]{geometry}

% -- Color y gráficos -----------------------------------------
\usepackage{xcolor}
\usepackage{graphicx}
\usepackage{tikz}
\usetikzlibrary{shapes.geometric, arrows.meta, positioning, fit, backgrounds, calc}

% -- Código fuente --------------------------------------------
\usepackage{listings}
\usepackage{inconsolata}   % fuente monoespaciada

% -- Tablas ---------------------------------------------------
\usepackage{booktabs}
\usepackage{tabularx}
\usepackage{multirow}
\usepackage{array}
\usepackage{colortbl}

% -- Otros ----------------------------------------------------
\usepackage{hyperref}
\usepackage{enumitem}
\usepackage{fancyhdr}
\usepackage{titlesec}
\usepackage{mdframed}
\usepackage{fontawesome5}
\usepackage{tcolorbox}
\tcbuselibrary{listings, skins, breakable}

% ============================================================
%  PALETA DE COLORES
% ============================================================
\definecolor{verdeAgro}   {HTML}{1D9E75}
\definecolor{verdeCl}     {HTML}{E1F5EE}
\definecolor{azulAPI}     {HTML}{185FA5}
\definecolor{azulCl}      {HTML}{E6F1FB}
\definecolor{ambarDB}     {HTML}{BA7517}
\definecolor{ambarCl}     {HTML}{FAEEDA}
\definecolor{rojoDel}     {HTML}{A32D2D}
\definecolor{rojoCl}      {HTML}{FCEBEB}
\definecolor{grisOscuro}  {HTML}{2C2C2A}
\definecolor{grisMedio}   {HTML}{5F5E5A}
\definecolor{grisCl}      {HTML}{F1EFE8}
\definecolor{fondoCodigo} {HTML}{1A1A18}
\definecolor{verdeCodigo} {HTML}{A0E8C0}
\definecolor{azulCodigo}  {HTML}{85B7EB}
\definecolor{naranjaKw}   {HTML}{EF9F27}

% ============================================================
%  ESTILO DE CÓDIGO -- JavaScript / JSON / Bash
% ============================================================
\lstdefinestyle{jsStyle}{
  backgroundcolor = \color{fondoCodigo},
  basicstyle      = \ttfamily\footnotesize\color{verdeCodigo},
  keywordstyle    = \color{naranjaKw}\bfseries,
  commentstyle    = \color{grisMedio}\itshape,
  stringstyle     = \color{azulCodigo},
  numberstyle     = \tiny\color{grisMedio},
  numbers         = left,
  numbersep       = 8pt,
  stepnumber      = 1,
  frame           = none,
  breaklines      = true,
  breakatwhitespace= false,
  breakindent     = 10pt,
  postbreak       = \mbox{\textcolor{grisMedio}{$\hookrightarrow$}\space},
  tabsize         = 2,
  showstringspaces= false,
  columns         = flexible,
  keepspaces      = true,
  xleftmargin     = 14pt,
  xrightmargin    = 4pt,
  aboveskip       = 4pt,
  belowskip       = 4pt,
}
\lstdefinestyle{jsonStyle}{
  style           = jsStyle,
  morestring      = [b]",
  literate        = {:}{{\textcolor{naranjaKw}{:}}}{1}
                    {,}{{\textcolor{grisMedio}{,}}}{1},
}
\lstdefinestyle{bashStyle}{
  style           = jsStyle,
  keywordstyle    = \color{verdeAgro}\bfseries,
  morekeywords    = {curl, node, npm, mkdir, cd},
}
\lstdefinestyle{xmlStyle}{
  style           = jsStyle,
  keywordstyle    = \color{azulCodigo}\bfseries,
  morestring      = [b]",
  morecomment     = [s]{<!--}{-->},
}

% ============================================================
%  CAJAS PERSONALIZADAS (tcolorbox)
% ============================================================
\tcbset{
  notaStyle/.style = {
    colback    = ambarCl,
    colframe   = ambarDB,
    fonttitle  = \bfseries\small,
    title      = {\faInfoCircle\ Nota},
    arc        = 4pt,
    boxrule    = 0.6pt,
    breakable,
  },
  codigoBox/.style = {
    colback    = fondoCodigo,
    colframe   = grisMedio,
    arc        = 6pt,
    boxrule    = 0.4pt,
    top        = 4pt, bottom = 4pt,
    left       = 6pt, right  = 6pt,
    breakable,
  },
  pasoStyle/.style = {
    colback    = verdeCl,
    colframe   = verdeAgro,
    fonttitle  = \bfseries\color{verdeAgro},
    arc        = 6pt,
    boxrule    = 0.8pt,
    breakable,
  },
  peligroStyle/.style = {
    colback    = rojoCl,
    colframe   = rojoDel,
    fonttitle  = \bfseries\small,
    title      = {\faExclamationTriangle\ Importante},
    arc        = 4pt,
    boxrule    = 0.6pt,
    breakable,
  },
}

% ============================================================
%  ENCABEZADO Y PIE DE PÁGINA
% ============================================================
\pagestyle{fancy}
\fancyhf{}
\fancyhead[L]{\small\color{grisMedio}\textit{AgroAPI Ecuador -- Tutorial Node.js + REST}}
\fancyhead[R]{\small\color{grisMedio}\textit{\nouppercase{\leftmark}}}
\fancyfoot[C]{\small\color{grisMedio}\thepage}
\renewcommand{\headrulewidth}{0.4pt}

% ============================================================
%  FORMATO DE SECCIONES
% ============================================================
\titleformat{\section}
  {\Large\bfseries\color{verdeAgro}}
  {\thesection.}{0.6em}{}[\vspace{-4pt}\rule{\linewidth}{0.6pt}]
\titleformat{\subsection}
  {\large\bfseries\color{azulAPI}}
  {\thesubsection.}{0.5em}{}
\titleformat{\subsubsection}
  {\normalsize\bfseries\color{grisOscuro}}
  {\thesubsubsection.}{0.4em}{}

% ============================================================
%  HIPERENLACES
% ============================================================
\hypersetup{
  colorlinks = true,
  linkcolor  = azulAPI,
  urlcolor   = verdeAgro,
  citecolor  = grisMedio,
  pdftitle   = {Tutorial API REST con Node.js -- Gestión de Cultivos},
  pdfauthor  = {AgroAPI Ecuador},
}

% ============================================================
%  COLORES MÉTODO HTTP
% ============================================================
\newcommand{\httpGET}  {\colorbox{azulCl}{\color{azulAPI}\textbf{\ttfamily GET}}}
\newcommand{\httpPOST} {\colorbox{verdeCl}{\color{verdeAgro}\textbf{\ttfamily POST}}}
\newcommand{\httpPUT}  {\colorbox{ambarCl}{\color{ambarDB}\textbf{\ttfamily PUT}}}
\newcommand{\httpDEL}  {\colorbox{rojoCl}{\color{rojoDel}\textbf{\ttfamily DEL}}}

% ============================================================
%  DOCUMENTO
% ============================================================
\begin{document}

% -- PORTADA -------------------------------------------------
\begin{titlepage}
  \pagecolor{verdeAgro}
  \color{white}
  \vspace*{3cm}
  \begin{center}
    {\fontsize{48}{52}\selectfont \bfseries \faLeaf}\\[1em]
    {\fontsize{28}{32}\selectfont \bfseries AgroAPI Ecuador}\\[0.5em]
    {\Large Tutorial Completo de API REST}\\[0.4em]
    {\large con Node.js, Express y JavaScript}\\[3cm]
    \begin{tikzpicture}
      \draw[white, line width=1.5pt] (0,0) -- (10,0);
    \end{tikzpicture}\\[1cm]
    {\large \textbf{Gestión de Cultivos Agrícolas}}\\[0.5em]
    {\normalsize Para principiantes absolutos en APIs web}\\[3cm]
    {\small Guayaquil, Ecuador \textbullet\ 2025}
  \end{center}
  \vfill
  \begin{center}
    {\small \faCode\ JSON \textbullet\ XML \textbullet\ WebServices \textbullet\ fetch() \textbullet\ PostgreSQL}
  \end{center}
\end{titlepage}
\pagecolor{white}\color{grisOscuro}

% -- ÍNDICE --------------------------------------------------
\tableofcontents
\newpage

% ============================================================
\section{¿Qué es una API y por qué la necesitamos?}
% ============================================================

Una \textbf{API} (Application Programming Interface, o Interfaz de Programación de Aplicaciones) es un conjunto de reglas y protocolos que permite que dos programas se comuniquen entre sí. En el contexto web, una \textbf{API REST} es un servidor que expone datos a través de URLs y responde con formato \texttt{JSON}.

\begin{tcolorbox}[notaStyle]
  Piensa en una API como un mesero en un restaurante: tú (el cliente) le dices qué quieres (petición HTTP), el mesero va a la cocina (servidor), y te trae la respuesta (datos JSON). No necesitas saber cómo funciona la cocina.
\end{tcolorbox}

\subsection{Arquitectura del sistema}

El sistema completo tiene tres capas que se comunican entre sí:

\bigskip

% -- DIAGRAMA TikZ: Arquitectura de 3 capas ------------------
\begin{center}
\resizebox{0.96\linewidth}{!}{%
\begin{tikzpicture}[
  capa/.style  = {draw, rounded corners=6pt,
                  minimum width=13.2cm, minimum height=2.6cm,
                  inner sep=8pt},
  nodo/.style  = {draw, rounded corners=4pt, fill=white,
                  minimum width=3.0cm, minimum height=1.2cm,
                  align=center, font=\small\strut},
  arr/.style   = {-Stealth, line width=1pt},
  lbl/.style   = {font=\footnotesize, align=center},
]
\pgfmathsetmacro{\W}{14.0}   % ancho total en cm (se ajusta con scalebox)

% -- CAPA CLIENTE -----------------------------------------
\node[capa, fill=azulCl, draw=azulAPI!60, text=azulAPI]
      at (0,0) (CLIENTE) {};
\node[font=\small\bfseries\color{azulAPI}]
      at (0, 0.85) {\faDesktop\; Capa Cliente (Navegador)};
\node[nodo, draw=azulAPI!40, fill=azulCl!40] at (-3.5,-0.1) {HTML + CSS\\{\tiny Interfaz}};
\node[nodo, draw=azulAPI!40, fill=azulCl!40] at ( 0.0,-0.1) {JavaScript\\{\tiny Lógica}};
\node[nodo, draw=azulAPI!40, fill=azulCl!40] at ( 3.5,-0.1) {\texttt{fetch()}\\{\tiny Peticiones HTTP}};

% -- FLECHAS CLIENTE ↔ SERVIDOR ---------------------------
\draw[arr, color=verdeAgro]
  (-1.6,-1.35) -- (-1.6,-1.95)
  node[midway, left=2pt, lbl, color=verdeAgro]
       {\faArrowDown\; Request};
\draw[arr, color=rojoDel]
  (1.6,-1.95) -- (1.6,-1.35)
  node[midway, right=2pt, lbl, color=rojoDel]
       {\faArrowUp\; Response JSON};

% -- CAPA SERVIDOR ----------------------------------------
\node[capa, fill=verdeCl, draw=verdeAgro!60]
      at (0,-3.3) (SERVIDOR) {};
\node[font=\small\bfseries\color{verdeAgro}]
      at (0,-2.45) {\faServer\; Servidor -- Node.js + Express (API REST)};
\node[nodo, draw=verdeAgro!40, fill=verdeCl!40] at (-3.5,-3.4)
      {Router\\{\tiny /api/cultivos}\\{\tiny /api/parcelas}};
\node[nodo, draw=verdeAgro!40, fill=verdeCl!40] at ( 0.0,-3.4)
      {Controladores\\{\tiny Lógica de negocio}\\{\tiny Validaciones}};
\node[nodo, draw=verdeAgro!40, fill=verdeCl!40] at ( 3.5,-3.4)
      {Middleware\\{\tiny cors(), json()}\\{\tiny Errores}};

% -- FLECHAS SERVIDOR ↔ DATOS -----------------------------
\draw[arr, color=verdeAgro]
  (-1.6,-4.65) -- (-1.6,-5.25)
  node[midway, left=2pt, lbl, color=verdeAgro]
       {\faArrowDown\; Query};
\draw[arr, color=ambarDB]
  (1.6,-5.25) -- (1.6,-4.65)
  node[midway, right=2pt, lbl, color=ambarDB]
       {\faArrowUp\; Datos};

% -- CAPA DATOS -------------------------------------------
\node[capa, fill=ambarCl, draw=ambarDB!60]
      at (0,-6.6) (DATOS) {};
\node[font=\small\bfseries\color{ambarDB}]
      at (0,-5.75) {\faDatabase\; Capa de Datos};
\node[nodo, draw=ambarDB!40, fill=ambarCl!40] at (-3.5,-6.7)
      {\texttt{cultivos.json}\\{\tiny Base de datos JSON}};
\node[nodo, draw=ambarDB!40, fill=ambarCl!40] at ( 0.0,-6.7)
      {\texttt{cultivos.xml}\\{\tiny Base de datos XML}};
\node[nodo, draw=grisMedio!40, fill=grisCl]   at ( 3.5,-6.7)
      {PostgreSQL\\{\tiny Producción}};

\end{tikzpicture}%
}
\end{center}

\bigskip

\subsection{Verbos HTTP y su significado}

Las operaciones CRUD se mapean directamente a verbos HTTP:

\begin{center}
\begin{tabular}{>{\bfseries}p{2.2cm} p{3cm} p{4cm} p{4cm}}
  \toprule
  \rowcolor{grisCl}
  Verbo HTTP & Operación CRUD & Qué hace & Ejemplo \\
  \midrule
  \httpGET   & \textbf{Read}   (Leer)     & Obtener datos       & \texttt{GET /api/cultivos} \\
  \httpPOST  & \textbf{Create} (Crear)    & Insertar nuevo dato & \texttt{POST /api/cultivos} \\
  \httpPUT   & \textbf{Update} (Actualizar) & Modificar un dato & \texttt{PUT /api/cultivos/3} \\
  \httpDEL   & \textbf{Delete} (Eliminar) & Borrar un dato      & \texttt{DELETE /api/cultivos/3} \\
  \bottomrule
\end{tabular}
\end{center}

% ============================================================
\section{Estructura del Proyecto}
% ============================================================

\begin{tcolorbox}[pasoStyle, title={\faFolderOpen\ Paso 1: Crear las carpetas}]
Ejecuta estos comandos en tu terminal:
\begin{lstlisting}[style=bashStyle]
mkdir cultivos-api
cd cultivos-api
mkdir data routes cliente
\end{lstlisting}
\end{tcolorbox}

La estructura final del proyecto es:

\begin{center}
\resizebox{\linewidth}{!}{%
\begin{tikzpicture}[
  raiz/.style  = {draw, rounded corners=4pt, fill=ambarCl,  draw=ambarDB,
                  font=\small\bfseries\ttfamily, minimum width=3.8cm,
                  minimum height=0.7cm, align=center},
  carpeta/.style={draw, rounded corners=4pt, fill=verdeCl, draw=verdeAgro,
                  font=\small\ttfamily, minimum width=3.2cm,
                  minimum height=0.65cm, align=center},
  archivo/.style={draw, rounded corners=3pt, fill=azulCl,  draw=azulAPI!60,
                  font=\footnotesize\ttfamily, minimum width=3.2cm,
                  minimum height=0.6cm,  align=center},
  raizarch/.style={draw, rounded corners=3pt, fill=grisCl, draw=grisMedio!60,
                  font=\footnotesize\ttfamily, minimum width=3.2cm,
                  minimum height=0.6cm,  align=center},
  lin/.style   = {draw=grisMedio!60, line width=0.5pt},
]

% Raíz
\node[raiz] (raiz) at (0,0) {cultivos-api/};

% -- Nivel 1: carpetas y archivos raíz en una fila ---------
% Columna data/
\node[carpeta] (data)   at (-5.0, -1.5) {data/};
% Columna routes/
\node[carpeta] (routes) at (-1.5, -1.5) {routes/};
% Columna cliente/
\node[carpeta] (cli)    at ( 1.5, -1.5) {cliente/};
% Archivos raíz
\node[raizarch] (srv)   at ( 4.2, -1.5) {server.js};
\node[raizarch] (pkg)   at ( 4.2, -2.4) {package.json};

% -- Nivel 2: hijos de cada carpeta ------------------------
\node[archivo] (json)  at (-5.8, -3.0) {cultivos.json};
\node[archivo] (xml)   at (-4.2, -3.0) {cultivos.xml};

\node[archivo] (cr)    at (-2.2, -3.0) {cultivosRoutes.js};
\node[archivo] (pr)    at ( 0.0, -3.0) {parcelasRoutes.js};

\node[archivo] (html)  at ( 0.8, -3.0) {index.html};
\node[archivo] (appjs) at ( 2.4, -3.0) {app.js};

% -- Líneas de conexión ------------------------------------
% raíz a nivel 1
\draw[lin] (raiz.south) -- ++(0,-0.4) -| (data.north);
\draw[lin] (raiz.south) -- ++(0,-0.4) -| (routes.north);
\draw[lin] (raiz.south) -- ++(0,-0.4) -| (cli.north);
\draw[lin] (raiz.south) -- ++(0,-0.4) -| (srv.north);
\draw[lin] (srv.south)  --             (pkg.north);

% data a sus hijos
\draw[lin] (data.south) -- ++(0,-0.3) -| (json.north);
\draw[lin] (data.south) -- ++(0,-0.3) -| (xml.north);

% routes a sus hijos
\draw[lin] (routes.south) -- ++(0,-0.3) -| (cr.north);
\draw[lin] (routes.south) -- ++(0,-0.3) -| (pr.north);

% cliente a sus hijos
\draw[lin] (cli.south) -- ++(0,-0.3) -| (html.north);
\draw[lin] (cli.south) -- ++(0,-0.3) -| (appjs.north);

\end{tikzpicture}%
}
\end{center}

% ============================================================
\section{Bases de Datos: JSON y XML}
% ============================================================

\subsection{¿Por qué JSON y XML?}

Para aprender, usamos archivos en lugar de una base de datos real. Ambos formatos pueden ser leídos directamente por Node.js sin instalar nada extra:

\begin{center}
\begin{tabular}{p{3cm} p{5cm} p{5cm}}
  \toprule
  \rowcolor{grisCl}
                & \textbf{JSON} & \textbf{XML} \\
  \midrule
  Sintaxis       & Llaves y corchetes \texttt{\{\}} & Etiquetas \texttt{<tag>} \\
  Nativo en JS   & \faCheck\ Sí (\texttt{JSON.parse()}) & \faTimesCircle\ Necesita \texttt{xml2js} \\
  Peso           & Más ligero & Más verboso \\
  Uso actual     & APIs REST modernas & Legado, SOAP, configs \\
  \bottomrule
\end{tabular}
\end{center}

\subsection{Archivo \texttt{data/cultivos.json}}

\begin{tcolorbox}[codigoBox]
\begin{lstlisting}[style=jsonStyle, language={}]
{
  "cultivos": [
    {
      "id": 1,
      "nombre": "Maiz Amarillo",
      "variedad": "Hibrido DK-7088",
      "fecha_siembra": "2024-11-15",
      "fecha_cosecha_estimada": "2025-04-20",
      "estado": "en_crecimiento",
      "parcela_id": 1,
      "area_hectareas": 3.5,
      "rendimiento_esperado_qq": 140,
      "observaciones": "Buen desarrollo vegetativo"
    },
    {
      "id": 2,
      "nombre": "Arroz",
      "variedad": "INIA Sequia",
      "fecha_siembra": "2024-12-01",
      "estado": "germinacion",
      "parcela_id": 2,
      "area_hectareas": 5.0
    }
  ],
  "parcelas": [
    {
      "id": 1,
      "nombre": "Parcela Norte A",
      "ubicacion": "Km 24 via Daule",
      "area_total_ha": 5.0,
      "tipo_suelo": "Franco arcilloso"
    }
  ],
  "insumos": [
    {
      "id": 1,
      "nombre": "Urea 46%",
      "tipo": "fertilizante",
      "unidad": "kg",
      "stock": 500,
      "precio_unitario": 0.75
    }
  ]
}
\end{lstlisting}
\end{tcolorbox}

\subsection{Archivo \texttt{data/cultivos.xml}}

\begin{tcolorbox}[codigoBox]
\begin{lstlisting}[style=xmlStyle, language={}]
<?xml version="1.0" encoding="UTF-8"?>
<agro_sistema>

  <cultivos>
    <cultivo id="1">
      <nombre>Maiz Amarillo</nombre>
      <variedad>Hibrido DK-7088</variedad>
      <fecha_siembra>2024-11-15</fecha_siembra>
      <estado>en_crecimiento</estado>
      <parcela_id>1</parcela_id>
      <area_hectareas>3.5</area_hectareas>
    </cultivo>
    <cultivo id="2">
      <nombre>Arroz</nombre>
      <variedad>INIA Sequia</variedad>
      <fecha_siembra>2024-12-01</fecha_siembra>
      <estado>germinacion</estado>
      <parcela_id>2</parcela_id>
      <area_hectareas>5.0</area_hectareas>
    </cultivo>
  </cultivos>

  <parcelas>
    <parcela id="1">
      <nombre>Parcela Norte A</nombre>
      <ubicacion>Km 24 via Daule</ubicacion>
      <area_total_ha>5.0</area_total_ha>
      <tipo_suelo>Franco arcilloso</tipo_suelo>
    </parcela>
  </parcelas>

</agro_sistema>
\end{lstlisting}
\end{tcolorbox}

% ============================================================
\section{El Servidor: Node.js + Express}
% ============================================================

\begin{tcolorbox}[pasoStyle, title={\faNpm\ Paso 2: Instalar dependencias}]
\begin{lstlisting}[style=bashStyle]
# Crear el package.json
npm init -y

# Instalar Express (servidor web) y CORS (seguridad)
npm install express cors xml2js

# Instalar nodemon para desarrollo (recarga automatica)
npm install --save-dev nodemon
\end{lstlisting}
\end{tcolorbox}

\subsection{Flujo de una petición HTTP}

\bigskip
\begin{center}
\resizebox{\linewidth}{!}{%
\begin{tikzpicture}[
  paso/.style  = {draw, rounded corners=5pt, minimum width=3.6cm,
                  minimum height=0.9cm, align=center, font=\small},
  arr/.style   = {-Stealth, line width=1.2pt},
  num/.style   = {circle, fill=verdeAgro, text=white,
                  font=\bfseries\scriptsize, inner sep=2.5pt},
]
% -- Nodos izquierda (petición) --------------------------
\node[paso, fill=azulCl,   draw=azulAPI!60]   (nav)  at (0,0)    {\faDesktop\\ Navegador};
\node[paso, fill=verdeCl,  draw=verdeAgro!60] (mid)  at (3.5,0)  {\faCogs\\ Middleware};
\node[paso, fill=verdeCl,  draw=verdeAgro!60] (rtr)  at (7.0,0)  {\faRoute\\ Router};
\node[paso, fill=verdeCl,  draw=verdeAgro!60] (ctrl) at (10.5,0) {\faCode\\ Controlador};
\node[paso, fill=ambarCl,  draw=ambarDB!60]   (db)   at (10.5,-2.2) {\faDatabase\\ cultivos.json};

% -- Flechas petición (izq -> der) ---------------------
\draw[arr, verdeAgro] (nav)  -- node[above,font=\tiny]{GET /api/cultivos} (mid);
\draw[arr, verdeAgro] (mid)  -- (rtr);
\draw[arr, verdeAgro] (rtr)  -- (ctrl);
\draw[arr, verdeAgro] (ctrl) -- node[right,font=\tiny]{leer} (db);

% -- Flecha respuesta (curva inferior) ----------------
\draw[arr, rojoDel, bend right=20]
  (db.west) to node[below,font=\tiny, color=rojoDel]{HTTP 200 + JSON} (nav.east |- db.west);

% -- Numeración ---------------------------------------
\node[num] at (1.75, 0.45) {1};
\node[num] at (5.25, 0.45) {2};
\node[num] at (8.75, 0.45) {3};
\node[num] at (10.5,-1.1)  {4};
\node[num] at (5.25,-1.7)  {5};

\end{tikzpicture}%
}
\end{center}
\bigskip

\subsection{Archivo \texttt{server.js} -- Punto de entrada}

\begin{tcolorbox}[codigoBox]
\begin{lstlisting}[style=jsStyle, language=JavaScript, morekeywords={require,const,let}]
// server.js -- Punto de entrada del WebService
const express = require('express');
const cors    = require('cors');
const path    = require('path');

// -- PARA POSTGRESQL: descomenta estas lineas --------------
// const { Pool } = require('pg');
// const pool = new Pool({
//   host: 'localhost', port: 5432,
//   database: 'agro_db',
//   user: 'postgres', password: 'tu_password'
// });
// module.exports = { pool };
// ----------------------------------------------------

const app  = express();
const PORT = 3000;

// -- MIDDLEWARE --------------------------------------------
app.use(cors());           // Permite peticiones desde cualquier origen
app.use(express.json());   // Parsea el body JSON de POST/PUT
app.use(express.static(path.join(__dirname, 'cliente')));

// -- RUTAS -------------------------------------------------
app.use('/api/cultivos',  require('./routes/cultivosRoutes'));
app.use('/api/parcelas',  require('./routes/parcelasRoutes'));
app.use('/api/insumos',   require('./routes/insumosRoutes'));

// -- INFORMACION DE LA API ---------------------------------
app.get('/api', (req, res) => {
  res.json({
    mensaje: 'Bienvenido a AgroAPI Ecuador',
    version: '1.0.0',
    endpoints: {
      cultivos: 'GET|POST /api/cultivos',
      cultivo:  'GET|PUT|DELETE /api/cultivos/:id',
      parcelas: 'GET|POST /api/parcelas',
      insumos:  'GET|POST /api/insumos'
    }
  });
});

// -- MANEJO DE ERRORES -------------------------------------
app.use((req, res) => {
  res.status(404).json({ error: 'Ruta no encontrada' });
});

// -- INICIAR -----------------------------------------------
app.listen(PORT, () => {
  console.log(`Servidor en http://localhost:${PORT}`);
});
\end{lstlisting}
\end{tcolorbox}

% ============================================================
\section{Rutas CRUD: \texttt{cultivosRoutes.js}}
% ============================================================

\begin{tcolorbox}[codigoBox]
\begin{lstlisting}[style=jsStyle, language=JavaScript, morekeywords={require,const,let,async,await}]
// routes/cultivosRoutes.js
const express = require('express');
const router  = express.Router();
const fs      = require('fs');
const path    = require('path');

const DB_PATH = path.join(__dirname, '../data/cultivos.json');

// Funciones auxiliares para leer y escribir el archivo JSON
function leerDatos()     { return JSON.parse(fs.readFileSync(DB_PATH, 'utf-8')); }
function guardarDatos(d) { fs.writeFileSync(DB_PATH, JSON.stringify(d, null, 2)); }

// -- GET /api/cultivos -------------------------------------
// Devuelve todos los cultivos (con filtro opcional por estado)
router.get('/', (req, res) => {
  try {
    const datos = leerDatos();
    let cultivos = datos.cultivos;

    if (req.query.estado)
      cultivos = cultivos.filter(c => c.estado === req.query.estado);

    res.json({ total: cultivos.length, cultivos });

    // CON POSTGRESQL:
    // const r = await pool.query('SELECT * FROM cultivos');
    // res.json({ total: r.rowCount, cultivos: r.rows });

  } catch (e) { res.status(500).json({ error: e.message }); }
});

// -- GET /api/cultivos/:id ---------------------------------
router.get('/:id', (req, res) => {
  try {
    const id     = parseInt(req.params.id);
    const datos  = leerDatos();
    const cultivo = datos.cultivos.find(c => c.id === id);
    if (!cultivo)
      return res.status(404).json({ error: `Cultivo ${id} no encontrado` });
    res.json(cultivo);
  } catch (e) { res.status(500).json({ error: e.message }); }
});

// -- POST /api/cultivos ------------------------------------
router.post('/', (req, res) => {
  try {
    const { nombre, variedad, fecha_siembra, parcela_id } = req.body;
    if (!nombre || !variedad || !fecha_siembra || !parcela_id)
      return res.status(400).json({ error: 'Faltan campos obligatorios' });

    const datos    = leerDatos();
    const nuevoId  = Math.max(...datos.cultivos.map(c => c.id)) + 1;
    const nuevo    = { id: nuevoId, nombre, variedad, fecha_siembra,
                       estado: req.body.estado || 'germinacion',
                       parcela_id: parseInt(parcela_id),
                       area_hectareas: parseFloat(req.body.area_hectareas) || 0 };
    datos.cultivos.push(nuevo);
    guardarDatos(datos);
    res.status(201).json({ mensaje: 'Cultivo creado', cultivo: nuevo });

    // CON POSTGRESQL:
    // const r = await pool.query(
    //   'INSERT INTO cultivos(nombre,variedad,fecha_siembra,...) VALUES($1,$2,$3,...) RETURNING *',
    //   [nombre, variedad, fecha_siembra, ...]
    // );
    // res.status(201).json({ cultivo: r.rows[0] });

  } catch (e) { res.status(500).json({ error: e.message }); }
});

// -- PUT /api/cultivos/:id ---------------------------------
router.put('/:id', (req, res) => {
  try {
    const id    = parseInt(req.params.id);
    const datos = leerDatos();
    const idx   = datos.cultivos.findIndex(c => c.id === id);
    if (idx === -1)
      return res.status(404).json({ error: `Cultivo ${id} no encontrado` });

    datos.cultivos[idx] = { ...datos.cultivos[idx], ...req.body, id };
    guardarDatos(datos);
    res.json({ mensaje: 'Actualizado', cultivo: datos.cultivos[idx] });
  } catch (e) { res.status(500).json({ error: e.message }); }
});

// -- DELETE /api/cultivos/:id ------------------------------
router.delete('/:id', (req, res) => {
  try {
    const id    = parseInt(req.params.id);
    const datos = leerDatos();
    const idx   = datos.cultivos.findIndex(c => c.id === id);
    if (idx === -1)
      return res.status(404).json({ error: `Cultivo ${id} no encontrado` });

    const eliminado = datos.cultivos.splice(idx, 1)[0];
    guardarDatos(datos);
    res.json({ mensaje: 'Eliminado', cultivo_eliminado: eliminado });
  } catch (e) { res.status(500).json({ error: e.message }); }
});

module.exports = router;
\end{lstlisting}
\end{tcolorbox}

% ============================================================
\section{Tabla de Endpoints REST}
% ============================================================

\begin{center}
\begin{tabular}{>{\centering\arraybackslash}p{1.8cm} p{5cm} p{4cm} p{3.5cm}}
  \toprule
  \rowcolor{grisCl}
  \textbf{Método} & \textbf{URL} & \textbf{Descripción} & \textbf{Parámetros} \\
  \midrule
  \httpGET  & \texttt{/api}                    & Info de la API          & --- \\
  \httpGET  & \texttt{/api/cultivos}            & Listar todos            & \texttt{?estado=...} \\
  \httpGET  & \texttt{/api/cultivos/:id}        & Obtener uno             & \texttt{:id} número \\
  \httpPOST & \texttt{/api/cultivos}            & Crear nuevo             & Body JSON \\
  \httpPUT  & \texttt{/api/cultivos/:id}        & Actualizar              & \texttt{:id} + Body \\
  \httpDEL  & \texttt{/api/cultivos/:id}        & Eliminar                & \texttt{:id} número \\
  \midrule
  \httpGET  & \texttt{/api/parcelas}            & Listar parcelas         & --- \\
  \httpGET  & \texttt{/api/parcelas/:id}        & Obtener una             & \texttt{:id} número \\
  \httpPOST & \texttt{/api/parcelas}            & Crear parcela           & Body JSON \\
  \httpPUT  & \texttt{/api/parcelas/:id}        & Actualizar              & \texttt{:id} + Body \\
  \httpDEL  & \texttt{/api/parcelas/:id}        & Eliminar                & \texttt{:id} número \\
  \midrule
  \httpGET  & \texttt{/api/insumos}             & Listar insumos          & \texttt{?tipo=...} \\
  \httpPOST & \texttt{/api/insumos}             & Crear insumo            & Body JSON \\
  \httpPUT  & \texttt{/api/insumos/:id}         & Actualizar              & \texttt{:id} + Body \\
  \httpDEL  & \texttt{/api/insumos/:id}         & Eliminar                & \texttt{:id} número \\
  \bottomrule
\end{tabular}
\end{center}

% ============================================================
\section{Cómo Probar los WebServices}
% ============================================================

\subsection{Opción A: con \texttt{curl} (terminal)}

\begin{tcolorbox}[pasoStyle, title={\faTerminal\ Paso 3: Iniciar el servidor y probar}]
\begin{lstlisting}[style=bashStyle]
# Terminal 1: iniciar el servidor
node server.js

# Terminal 2: probar los endpoints
\end{lstlisting}
\end{tcolorbox}

\subsubsection{Peticiones GET}

\begin{tcolorbox}[codigoBox]
\begin{lstlisting}[style=bashStyle]
# Listar todos los cultivos
curl http://localhost:3000/api/cultivos

# Filtrar por estado
curl "http://localhost:3000/api/cultivos?estado=en_crecimiento"

# Obtener el cultivo con ID 1
curl http://localhost:3000/api/cultivos/1

# Ver todas las parcelas
curl http://localhost:3000/api/parcelas
\end{lstlisting}
\end{tcolorbox}

\subsubsection{Petición POST -- Crear nuevo cultivo}

\begin{tcolorbox}[codigoBox]
\begin{lstlisting}[style=bashStyle]
curl -X POST http://localhost:3000/api/cultivos \
  -H "Content-Type: application/json" \
  -d '{
    "nombre": "Pimiento",
    "variedad": "California Wonder",
    "fecha_siembra": "2025-01-20",
    "parcela_id": 5,
    "area_hectareas": 0.5,
    "estado": "germinacion"
  }'
\end{lstlisting}
\end{tcolorbox}

\subsubsection{Petición PUT -- Actualizar estado}

\begin{tcolorbox}[codigoBox]
\begin{lstlisting}[style=bashStyle]
curl -X PUT http://localhost:3000/api/cultivos/2 \
  -H "Content-Type: application/json" \
  -d '{
    "estado": "en_crecimiento",
    "observaciones": "Malezas controladas, buen desarrollo"
  }'
\end{lstlisting}
\end{tcolorbox}

\subsubsection{Petición DELETE -- Eliminar}

\begin{tcolorbox}[codigoBox]
\begin{lstlisting}[style=bashStyle]
curl -X DELETE http://localhost:3000/api/cultivos/6
\end{lstlisting}
\end{tcolorbox}

\subsection{Opción B: con Postman o Thunder Client}

\begin{enumerate}[leftmargin=*, label=\textbf{\arabic*.}]
  \item Instala \textbf{Thunder Client} como extensión de VS Code, o descarga \href{https://www.postman.com}{\textbf{Postman}}.
  \item Crea una nueva petición.
  \item Selecciona el método (\texttt{GET}, \texttt{POST}, etc.) y escribe la URL.
  \item Para \texttt{POST} y \texttt{PUT}: ve a \textit{Body} -> \textit{JSON} y escribe el objeto.
  \item Haz clic en \textit{Send} y observa la respuesta.
\end{enumerate}

\begin{tcolorbox}[notaStyle]
  Los códigos de respuesta HTTP más importantes son: \textbf{200 OK} (éxito), \textbf{201 Created} (recurso creado), \textbf{400 Bad Request} (datos incorrectos), \textbf{404 Not Found} (no encontrado), \textbf{500 Server Error} (error del servidor).
\end{tcolorbox}

% ============================================================
\section{Consumir la API desde JavaScript}
% ============================================================

\subsection{La función \texttt{fetch()}}

\texttt{fetch()} es la función nativa de JavaScript en el navegador para hacer peticiones HTTP. Devuelve una \textbf{Promesa} (Promise), lo que significa que la respuesta llegará en el futuro.

\bigskip

\begin{center}
\resizebox{\linewidth}{!}{%
\begin{tikzpicture}[
  caja/.style  = {draw, rounded corners=5pt, minimum width=3.4cm,
                  minimum height=0.9cm, align=center, font=\small},
  arr/.style   = {-Stealth, line width=1pt, color=verdeAgro},
  lbl/.style   = {font=\footnotesize\ttfamily, color=grisMedio},
]
% -- Nodos en fila horizontal --------------------------
\node[caja, fill=azulCl, draw=azulAPI!60]    (A) at (0.0, 0) {\texttt{fetch(url)}\\{\tiny llama a la API}};
\node[caja, fill=ambarCl,draw=ambarDB!60]    (B) at (4.0, 0) {Promise\\{\tiny respuesta futura}};
\node[caja, fill=verdeCl, draw=verdeAgro!60] (C) at (8.0, 0) {\texttt{res.json()}\\{\tiny convierte texto}};
\node[caja, fill=verdeCl, draw=verdeAgro!60] (D) at (8.0,-2.0) {Datos JS\\{\tiny objeto/array}};
\node[caja, fill=azulCl,  draw=azulAPI!60]   (E) at (4.0,-2.0) {Actualizar DOM\\{\tiny mostrar HTML}};

% -- Flechas -------------------------------------------
\draw[arr] (A) -- node[above, lbl]{.then(res =>)} (B);
\draw[arr] (B) -- node[above, lbl]{res =>}        (C);
\draw[arr] (C) -- node[right, lbl]{.then(data =>)}(D);
\draw[arr] (D) -- node[below, lbl]{data =>}       (E);

\end{tikzpicture}%
}
\end{center}

\bigskip

\subsection{Función central de consumo de API}

\begin{tcolorbox}[codigoBox]
\begin{lstlisting}[style=jsStyle, language=JavaScript, morekeywords={async,await,const,let}]
// cliente/app.js
// URL base de la API -- cambiar si el servidor esta en otro puerto
const BASE_URL = 'http://localhost:3000/api';

/**
 * Funcion central para hacer peticiones a la API.
 * url     : endpoint, ej: BASE_URL + '/cultivos'
 * metodo  : 'GET', 'POST', 'PUT', 'DELETE'
 * datos   : objeto JavaScript para POST/PUT (opcional)
 */
async function pedirAPI(url, metodo = 'GET', datos = null) {
  // 1. Configurar la peticion
  const opciones = {
    method:  metodo,
    headers: { 'Content-Type': 'application/json' }
  };

  // 2. Agregar body solo en POST y PUT
  if (datos && (metodo === 'POST' || metodo === 'PUT')) {
    opciones.body = JSON.stringify(datos);  // objeto JS -> texto JSON
  }

  // 3. Realizar la peticion HTTP
  const respuesta = await fetch(url, opciones);

  // 4. Convertir la respuesta de texto a objeto JavaScript
  const json = await respuesta.json();

  // 5. Si el servidor devolvio error (4xx, 5xx), lanzar excepcion
  if (!respuesta.ok) {
    throw new Error(json.error || `Error HTTP ${respuesta.status}`);
  }

  return json;  // devolver los datos listos para usar
}
\end{lstlisting}
\end{tcolorbox}

\subsection{Leer y mostrar cultivos}

\begin{tcolorbox}[codigoBox]
\begin{lstlisting}[style=jsStyle, language=JavaScript, morekeywords={async,await,const,let}]
// -- CARGAR todos los cultivos y mostrarlos en una tabla ---
async function cargarCultivos() {
  try {
    // Llamar a la API
    const data = await pedirAPI(`${BASE_URL}/cultivos`);

    // data.cultivos es un array de objetos
    const tbody = document.getElementById('tabla-cuerpo');

    // Construir las filas de la tabla con los datos recibidos
    tbody.innerHTML = data.cultivos.map(c => `
      <tr>
        <td>${c.id}</td>
        <td>${c.nombre}</td>
        <td>${c.variedad}</td>
        <td>${c.fecha_siembra}</td>
        <td>${c.area_hectareas} ha</td>
        <td>${c.estado}</td>
        <td>
          <button onclick="editarCultivo(${c.id})">Editar</button>
          <button onclick="eliminarCultivo(${c.id})">Eliminar</button>
        </td>
      </tr>
    `).join('');

  } catch (error) {
    console.error('Error al cargar cultivos:', error.message);
    alert('No se pudo conectar con la API');
  }
}

// Ejecutar al cargar la pagina
cargarCultivos();
\end{lstlisting}
\end{tcolorbox}

\subsection{Crear, editar y eliminar}

\begin{tcolorbox}[codigoBox]
\begin{lstlisting}[style=jsStyle, language=JavaScript, morekeywords={async,await,const,let}]
// -- CREAR un nuevo cultivo (POST) -------------------------
async function crearCultivo() {
  const nuevoCultivo = {
    nombre:        document.getElementById('campo-nombre').value,
    variedad:      document.getElementById('campo-variedad').value,
    fecha_siembra: document.getElementById('campo-siembra').value,
    parcela_id:    parseInt(document.getElementById('campo-parcela').value),
    area_hectareas:parseFloat(document.getElementById('campo-area').value),
    estado:        'germinacion'
  };

  try {
    const resultado = await pedirAPI(`${BASE_URL}/cultivos`, 'POST', nuevoCultivo);
    alert('Cultivo creado con ID: ' + resultado.cultivo.id);
    cargarCultivos();  // recargar la tabla
  } catch (e) {
    alert('Error: ' + e.message);
  }
}

// -- ACTUALIZAR un cultivo existente (PUT) -----------------
async function actualizarEstado(id, nuevoEstado) {
  try {
    await pedirAPI(`${BASE_URL}/cultivos/${id}`, 'PUT', { estado: nuevoEstado });
    cargarCultivos();
  } catch (e) {
    alert('Error al actualizar: ' + e.message);
  }
}

// -- ELIMINAR un cultivo (DELETE) --------------------------
async function eliminarCultivo(id) {
  if (!confirm(`Eliminar cultivo #${id}?`)) return;
  try {
    await pedirAPI(`${BASE_URL}/cultivos/${id}`, 'DELETE');
    cargarCultivos();  // actualizar la tabla
  } catch (e) {
    alert('Error al eliminar: ' + e.message);
  }
}
\end{lstlisting}
\end{tcolorbox}

% ============================================================
\section{Migración a PostgreSQL}
% ============================================================

Cuando el proyecto crezca, puedes reemplazar los archivos JSON por una base de datos real. Los cambios necesarios son mínimos:

\begin{center}
\begin{tabular}{p{6cm} p{6.5cm}}
  \toprule
  \rowcolor{grisCl}
  \textbf{Con JSON (ahora)} & \textbf{Con PostgreSQL (producción)} \\
  \midrule
  \texttt{fs.readFileSync(DB\_PATH)} & \texttt{await pool.query('SELECT * FROM...')} \\
  \texttt{JSON.parse(contenido)}     & Resultados en \texttt{resultado.rows} \\
  \texttt{fs.writeFileSync(...)}     & \texttt{await pool.query('INSERT INTO...')} \\
  Autogenerar ID con \texttt{Math.max} & \texttt{SERIAL} o \texttt{GENERATED ALWAYS AS IDENTITY} \\
  \bottomrule
\end{tabular}
\end{center}

\begin{tcolorbox}[codigoBox]
\begin{lstlisting}[style=jsStyle, language=JavaScript, morekeywords={const,require,await}]
// Ejemplo: GET /api/cultivos con PostgreSQL
const { Pool } = require('pg');
const pool = new Pool({
  host: 'localhost', port: 5432,
  database: 'agro_db',
  user: 'postgres', password: 'mi_password'
});

router.get('/', async (req, res) => {
  try {
    const resultado = await pool.query(
      'SELECT * FROM cultivos ORDER BY id'
    );
    res.json({ total: resultado.rowCount, cultivos: resultado.rows });
  } catch (e) { res.status(500).json({ error: e.message }); }
});
\end{lstlisting}
\end{tcolorbox}

% ============================================================
\section{Resumen y Próximos Pasos}
% ============================================================

\subsection{Lo que aprendiste en este tutorial}

\begin{itemize}[leftmargin=*, label={\color{verdeAgro}\faCheckCircle}]
  \item \textbf{Arquitectura REST}: separación en capas Cliente $\rightarrow$ API $\rightarrow$ Datos.
  \item \textbf{Verbos HTTP}: GET (leer), POST (crear), PUT (actualizar), DELETE (eliminar).
  \item \textbf{Bases de datos de archivo}: JSON y XML como almacenamiento simple.
  \item \textbf{Node.js + Express}: crear un servidor HTTP con rutas parametrizadas.
  \item \textbf{Middleware}: \texttt{cors()}, \texttt{express.json()}.
  \item \textbf{fetch() en el cliente}: peticiones asíncronas con \texttt{async/await}.
  \item \textbf{Manejo de errores}: códigos HTTP y bloques \texttt{try/catch}.
  \item \textbf{Migración a PostgreSQL}: qué líneas cambiar cuando el proyecto crezca.
\end{itemize}

\subsection{Próximos pasos sugeridos}

\begin{enumerate}[leftmargin=*, label=\textbf{\arabic*.}]
  \item \textbf{Autenticación}: agrega JWT (\texttt{jsonwebtoken}) para proteger los endpoints.
  \item \textbf{Validación}: usa \texttt{joi} o \texttt{express-validator} para validar los datos del body.
  \item \textbf{Base de datos real}: migra a PostgreSQL con \texttt{pg} o usa un ORM como \texttt{Sequelize} o \texttt{Prisma}.
  \item \textbf{Variables de entorno}: mueve las credenciales a un archivo \texttt{.env} con \texttt{dotenv}.
  \item \textbf{Despliegue}: publica tu API en \textbf{Railway}, \textbf{Render} o \textbf{Heroku}.
\end{enumerate}

\begin{tcolorbox}[peligroStyle]
  Nunca subas contraseñas, claves API o credenciales de base de datos a GitHub. Siempre usa archivos \texttt{.env} y agrégalos al \texttt{.gitignore}.
\end{tcolorbox}

\vfill
\begin{center}
  \color{grisMedio}
  \rule{10cm}{0.4pt}\\[0.5em]
  {\small \faLeaf\ AgroAPI Ecuador $\cdot$ Tutorial Node.js + REST $\cdot$ 2025}\\
  {\small Compilado con \LaTeX\ $\cdot$ \texttt{pdflatex tutorial\_api.tex}}
\end{center}

\end{document}